\documentclass{beamer}

% Selección del tema y colores
\usetheme{Madrid}
\usecolortheme{default}

% Paquetes de idioma y matemáticas
\usepackage[utf8]{inputenc}
\usepackage[spanish]{babel}
\usepackage{amsmath}
\usepackage{amsfonts}
\usepackage{amssymb}
\usepackage{mathtools} % Necesario para \coloneqq

% Información de la portada
\title[Método Adomiano (ADM)]{El Método de Descomposición Adomiano}
\subtitle{Estructura y Aplicaciones}
\author[Matihus Molina]{Matihus Alberth Molina Larios}
\institute[UNMSM]{
	Universidad Nacional Mayor de San Marcos \\
	Facultad de Ciencias Matemáticas / Matemática 
}
\date[Abril 2026]{17 de abril de 2026 \\ San Martín de Porres, Lima, Perú}

% --- NUEVO PAQUETE AÑADIDO AQUÍ ---
\usepackage{csquotes}

% Paquetes para bibliografía con archivo .bib
\usepackage[backend=biber, style=numeric, sorting=nty]{biblatex}
\addbibresource{references.bib} % Reemplaza con el nombre de tu archivo

\begin{document}
	
	% Diapositiva 1: Portada
	\begin{frame}
		\titlepage
	\end{frame}
	
	% Diapositiva 2: Índice
	\begin{frame}{Contenido}
		\tableofcontents
	\end{frame}
	
	% Sección 1: Introducción
	\section{Introducción}
	
	\begin{frame}{Introducción al ADM}
		\begin{itemize}
			\item Introducido y utilizado extensamente en la década de 1980.\cite{Abello_Res_2016}
			\item Ampliamente aplicado para aproximar las soluciones a diversas clases de ecuaciones diferenciales no lineales \textbf{sin recurrir a} métodos de discretización, perturbación, linealización o aproximaciones de cierre.\cite{Abello_Res_2016}
			\item Generalmente es necesario hacer suposiciones simplificadoras o imponer condiciones restrictivas, convirtiendo el problema en un Problema de Valor Inicial (PVI).\cite{Mak2021}
		\end{itemize}
	\end{frame}
	
	\begin{frame}{Naturaleza de los Problemas y Solución}
		\begin{itemize}
			\item Es útil para problemas de naturaleza \textbf{determinista y estocástica}, lineales y no lineales, provenientes de la física, química, biología e ingeniería.\cite{Abello_Res_2016}
			\item La solución se obtiene mediante una expansión en serie de los llamados \textbf{polinomios de Adomian}.
			\item El método proporciona una serie de rápida convergencia cuyos componentes se calculan de forma elegante e individual, paso a paso, mediante una relación recursiva.\cite{Wazwaz2001}
			\item \textit{Desafío actual:} Una pregunta más difícil de responder es cómo este método puede ser modificado para abordar el concepto de puntos singulares.\cite{Wazwaz2001}
		\end{itemize}
	\end{frame}
	
	% Sección 2: Fundamento Matemático
	\section{Fundamento Matemático}
	
	\begin{frame}{Estructura de la Ecuación}
		La ecuación en forma de operador general se escribe como:
		\begin{equation}\label{eq:generaloperator}
			Lu+Ru+Nu=g.
		\end{equation}
		Donde los operadores representan:
		\begin{itemize}
			\setlength\itemsep{0em}
			\item $L$: Derivada de mayor orden.
			\item $R$: Derivadas de menor orden.
			\item $N$: Términos no lineales.
			\item $g$: Término fuente.
		\end{itemize}
		
		Aplicando el operador inverso $L^{-1}$ a la ecuación \eqref{eq:generaloperator}, obtenemos:
		\begin{equation*}
			u = \Phi + L^{-1}g - L^{-1}Ru - L^{-1}Nu,
		\end{equation*}
		donde $\Phi$ se elige para que la solución $u$ satisfaga las condiciones iniciales.
	\end{frame}
	
	% Sección 3: Polinomios y Esquema Recursivo
	\section{Polinomios de Adomian y Esquema Recursivo}
	
	\begin{frame}{Polinomios de Adomian e Inversa Lane-Emden}
		\begin{definition}[Polinomios de Adomian $A_{n}$ y $L^{-1}$ para Lane-Emden]
			\begin{align*}
				\forall n\in\mathbb{N}\cup\left\{0\right\}:
				A_{n}  & \coloneqq
				\frac{1}{n!}
				\frac{d^n}{d\lambda^n}
				\left[
				N\left(
				\sum_{i=0}^{\infty}
				\lambda^{i}
				u_{i}
				\right)
				\right]_{\lambda=0} \\
				L^{-1} & \coloneqq
				\int_{0}^{\xi}
				\xi^{-2}
				\int_{0}^{\xi}
				\xi^{2}\left(\cdot\right)
				\,d\xi \,d\xi.
			\end{align*}
		\end{definition}
	\end{frame}
	
	\begin{frame}{Esquema Recursivo}
		Descomponiendo la solución como $u=\sum\limits_{n=0}^{\infty}u_{n}$, el esquema recursivo es:
		
		\vspace{0.5cm}
		\begin{equation}\label{eq:adm_system}
			\boxed{
				\left\{
				\begin{aligned}
					\forall n\in\mathbb{N}\cup\left\{0\right\}:
					u_{n+1} & = -L^{-1}Ru_{n}-L^{-1}A_{n}. \\
					u_{0}   & = \Phi+L^{-1}g.
				\end{aligned}
				\right.
			}
		\end{equation}
		\vspace{0.5cm}
		
		Este sistema permite determinar cada término de la serie analítica apoyándose en los polinomios calculados previamente.
	\end{frame}
	
	% Sección 4: Ejemplos
	\section{Ejemplos Resueltos}
	\begin{frame}{Ejemplos de Ecuaciones Resueltas por ADM}
		\begin{itemize}
			\item \textit{[Espacio reservado para tu lista de ejemplos de ecuaciones diferenciales]}
			\item Ecuación de Lane-Emden (aplicando el $L^{-1}$ definido anteriormente).
			\item ...
			\item ...
		\end{itemize}
	\end{frame}
	
	% Sección 5: Conclusión
	\section{Conclusiones}
	\begin{frame}{Conclusiones}
		\begin{itemize}
			\item El Método de Descomposición Adomiano es una herramienta poderosa y versátil para sistemas no lineales en múltiples disciplinas.
			\item Proporciona soluciones analíticas continuas y evita el enorme costo computacional de métodos puramente numéricos.
			\item Mantiene intacta la física del problema subyacente.
		\end{itemize}
		
		\vspace{1cm}
		\begin{center}
			\Large \textbf{¡Gracias por su atención!}
		\end{center}
	\end{frame}
	
	% Sección 6: Referencias
	\section{Referencias}
	\begin{frame}[allowframebreaks]{Referencias}
		% Esto imprimirá automáticamente todo lo que hayas citado de tu archivo .bib
		\printbibliography 
	\end{frame}
	
\end{document}